The general distinction between matching human response distributions and predicting their categories is established in FER and subjective language tasks \cite{barsoum2016ferplus,washington2021crowdsourcing,khurana2024crowd}. Our contribution is an empirical account of how specified target changes, output vocabulary and accepted images affect that distinction. The original Swin hard--soft difference falls from 2.61 pp to 0.38 pp against Uniform and does not persist in the ResNet comparison. The entropy-preserving permutation adds evidence that target probability values alone do not explain the observed ordering, although its changed class priors prevent a unique category-identity interpretation.

The eight-output results make the control dependence clearer. The distance-selected difference is 0.84 pp against Uniform but only 0.05 pp against Tie-Uniform; the latter changes direction under other confidence scores. Multiplying these own-ranking differences by 2,620 gives approximately 21.9 and 1.4 expected annotation-draw disagreements, respectively. These are count differences between separately selected populations, not the number of identical images corrected. The near-zero Tie-Uniform contrast does not prove equivalence, but it limits claims that detailed minority proportions generally improve selective decisions beyond preserving ties and dominant probabilities.

Identical-image comparisons add information that the own-ranking mean cannot supply. The predictor ordering changes with the selector's accepted set, so a lower selective risk does not establish uniformly better category choices on a fixed population. This pattern motivates reporting cross-evaluation alongside coverage curves; it does not identify a unique causal contribution of predictor versus selector. Likewise, the unsupported-mass decomposition is an algebraic identity. Its empirical role is to reveal how much of observed disagreement lies outside the available outputs and how that composition changes between accepted populations, not to establish a new theorem or diagnose internal emotion.

The follow-up target, vocabulary and ranking comparisons address explanations left open by the original results. Adding contempt changes supported weights and classifier dimension together and does not lower Soft's own-ranking risk in this comparison. The alternative stopping rule retains a Soft--Uniform advantage but does not resolve a consistent Soft--Tie-Uniform advantage. Lower squared distance across the matched conditions is partly expected from objective alignment; its value here is that it contrasts with the less stable categorical ordering, not that it establishes a new learning principle.

The external and calibration findings limit practical interpretation. Transferred empirical thresholds overshoot, and the conservative comparator sacrifices coverage. RAVDESS results are sensitive to checkpoint, seed and reference and exhibit substantial category-choice excess. Aggregation sensitivity can locate one pipeline dependence but cannot separate all effects of source population, cropping context and capture conditions. Confidence ranking is also a plausible explanation in light of label-smoothing results \cite{xia2025smoothing}; distributional distance alone cannot certify a useful rejector or calibrated risk \cite{baan2022calibration,guo2017calibration}. The supplementary preview consequently remains outside the validated operational scope.

\subsection{Limitations}
All training phases adaptively reuse FER2013 images with FER+ annotations. The final follow-up and locally frozen external plan preserve their timing records but do not restore an untouched static-image evaluation. Exact-pixel isolation and perceptual screening address particular overlap routes; identities and near-duplicate independence remain unverified. Conditional intervals describe resampling the observed pixels or actors for the retained fits. They omit uncertainty over raters, new seeds and populations and do not compensate for prior hypothesis selection.

Three paired seeds, two backbones and a fixed twelve-epoch recipe provide limited optimization evidence. The new eight-output stopping sensitivity tests an additional selection objective, not all schedules or hyperparameters. Vocabulary expansion and permutation address specified explanations rather than every effect of entropy, priors or regularization. Human vote frequencies are finite-panel observations, not emotion truth. RAVDESS is adult acted speech video with clip judgments, not a naturalistic messaging sample or independent static-image training replication.

No communication, demographic, cultural or mobile-efficiency study is supplied. Detector failures, pose, lighting and accessories may affect the local preview, whose update/freshness controls have not been validated through user research. Historical preprint manual-curation scores remain excluded because their records and checkpoints were not retained. These scope limits restrict generalization; they do not justify favourable unobserved results.
